\begin{abstract}
OpenClaw-like open-source agents have achieved broad deployment for multi-step tool-use and reasoning tasks, yet their failure modes remain poorly understood. Prior work either optimizes these agents with expensive training recipes or reports aggregate benchmark scores without diagnosing \emph{what} goes wrong and \emph{why}. We present a unified study that closes the loop from evaluation to diagnosis to repair. First, we build a pinned evaluation harness across two agent benchmarks (PinchBench and AgentBench-OpenClaw) and establish a systematic \textbf{failure taxonomy} with five dominant categories: (A) task understanding and planning drift, (B) tool and environment grounding failure, (C) memory and state management failure, (D) verification and recovery deficiency, and (E) skill utilization deficiency. Through failure case analysis on 4--5 consistently failing tasks per benchmark, we identify structural failure patterns: E-type failures (ineffective skill activation or composition) appear in 13--17\% of tasks and require different repair strategies depending on whether the skill is missing or simply underutilized; the memory domain on AgentBench-OC consistently scores 50.0 across all model checkpoints, indicating a structural C-type failure. Second, we design four classes of \textbf{lightweight training-free recipes} informed by the failure taxonomy, and provide a preliminary analysis of their expected effectiveness based on the observed failure patterns: E-type failures require either procedural knowledge injection (T4a) or skill-activation prompting (T4b), while A-type and B-type failures show potential for prompting-based repair. Third, we outline a controlled \textbf{small-data SFT study} (100--1000 samples) as the next phase to quantify complementary gains on category E. Our preliminary findings suggest that a large fraction of agent failures are structurally diagnosable and that the repair strategy must be matched to the failure type: skill-utilization failures require targeted interventions at the skill selection and activation level, while planning and grounding failures may be addressable by lightweight prompting recipes.
\end{abstract}
